\begin{figure}[!htbp]
\centering
\includegraphics[width=0.96\textwidth]{figures/ch2/seismic_vulnerability_rc_building.png}
\caption{Seismic vulnerability of low-rise RC buildings.}
\label{fig:rc_vulnerability_concept}
\end{figure}

Figure~\ref{fig:rc_vulnerability_concept} summarizes the main vulnerability mechanisms discussed in the literature for low-rise reinforced-concrete buildings. Inadequate detailing, workmanship and material quality can increase damage susceptibility, while plan irregularity and eccentricity can disturb stiffness and force distribution and increase torsional and drift concentration. The schematic uses the three-storey study context and an illustrative re-entrant plan based on the study grid only to communicate these mechanisms; it is not an analysis result or a fragility curve. This interpretation is consistent with the Nepalese RC-building observations and irregular-building studies reviewed in this chapter \cite{chaulagain2013seismic,parajuli2015gorkha,poudel2017seismic}.
\FloatBarrier
